\begin{table*}[t]
\centering
\caption{Main real-reuse experiment table. Filled scores come from local raw rows; future reruns or additional rows may be added. Original-paper scores are reported references unless the same data, input/output, metric, budget, and runtime setting are reproduced locally.}
\label{tab:real-reuse-main}
\resizebox{\textwidth}{!}{%
\begin{tabular}{llllllllll}
\toprule
Task & Paper & Domain & Input & Output & Metric & Reference & Summary & PaperToSkill & Status \\
\midrule
AIDE-T1 & AIDE & ML engineering & Kaggle-style dataset + metric & Runnable solution/submission & validation\_score & Reported AIDE ref. & 0.816 & 0.817 & Scored (GPT-family) \\
AIDE-T2 & AIDE & ML engineering & Weak ML script + feedback & Improved script + trajectory & best\_node\_score & Reported AIDE ref. & 0.000 & 0.826 & Scored (GPT-family) \\
SWE-T1 & SWE-agent & Software engineering & Repo issue + tests & Patch + test log & resolved & Reported SWE-agent ref. & 0.000 & 0.000 & Scored (GPT-family) \\
SWE-T2 & SWE-agent & Software engineering & Failing test + repo & Focused patch + verification & tests\_passed & Reported SWE-agent ref. & 0.000 & 1.000 & Scored (GPT-family) \\
REF-T1 & Reflexion & Reasoning / QA & Multi-hop QA + feedback & Final answer + reflection trace & exact\_match\_or\_f1 & Reported Reflexion ref. & 1.000 & 1.000 & Scored (GPT-family) \\
REF-T2 & Reflexion & Decision / programming & Failed attempt + checker feedback & Corrected second attempt & second\_attempt\_success & Reported Reflexion ref. & 1.000 & 1.000 & Scored (GPT-family) \\
SNAP-T1 & SnapATAC2 & Single-cell omics & Small single-cell dataset & Pipeline + embedding artifacts & runtime\_memory\_quality & Reported SnapATAC2 ref. & 0.000 & 0.500 & Scored (GPT-family) \\
SNAP-T2 & SnapATAC2 & Single-cell omics & Single-cell labels/proxy task & Clustering/marker artifacts & ari\_nmi\_runtime\_memory & Reported SnapATAC2 ref. & 0.200 & 0.400 & Scored (GPT-family) \\
\bottomrule
\end{tabular}%
}
\end{table*}

\begin{table*}[t]
\centering
\caption{Real-reuse failure-boundary analysis derived from the same scored raw rows as Table~\ref{tab:real-reuse-main}. This table explains first-pass boundary modes; it does not add new task-success evidence.}
\label{tab:real-reuse-failure-analysis}
\resizebox{\textwidth}{!}{%
\begin{tabular}{lllll}
\toprule
Task & Summary Outcome & PaperToSkill Outcome & Boundary Mode & Contract Implication \\
\midrule
AIDE-T1 & 0.816; success & 0.817; success & Solved by both & harder task slice or stricter scorer \\
AIDE-T2 & 0.000; timeout & 0.826; success & PaperToSkill-only success & preserve patch/tool-use constraints \\
SWE-T1 & 0.000; patch apply failed & 0.000; patch apply failed & Patch application & patch-format and apply-check contract \\
SWE-T2 & 0.000; patch apply failed & 1.000; success & PaperToSkill-only success & preserve patch/tool-use constraints \\
REF-T1 & 1.000; success & 1.000; success & Solved by both & harder task slice or stricter scorer \\
REF-T2 & 1.000; success & 1.000; success & Solved by both & harder task slice or stricter scorer \\
SNAP-T1 & 0.000; artifact parse failed & 0.500; artifact incomplete & Artifact completion & required artifact and metric manifest \\
SNAP-T2 & 0.200; artifact incomplete & 0.400; artifact incomplete & Artifact completion & required artifact and metric manifest \\
\bottomrule
\end{tabular}%
}
\end{table*}

\begin{table*}[t]
\centering
\caption{SWE-T1 shared-source-context follow-up. The first-pass SWE-T1 row remains the main-table evidence; this paired follow-up exposes the same locked SQLFluff source slice to both conditions and is diagnostic rather than a replacement.}
\label{tab:swe-t1-source-context-followup}
\resizebox{\textwidth}{!}{%
\begin{tabular}{llllllllllll}
\toprule
Task & Condition & First Run & First Score & First Failure & First Patch & Follow-up Run & Follow-up Score & Follow-up Patch & Follow-up Test & Follow-up Failure & Interpretation \\
\midrule
SWE-T1 & Summary & phase97\_gpt\_swe\_t1\_summary\_retry & 0.000 & patch\_apply\_failed & No & phase107\_gpt\_swe\_t1\_source\_context\_followup & 0.000 & Yes & No & test\_command\_failed & Source context fixed patch application; hidden test still failed \\
SWE-T1 & PaperToSkill & phase97\_gpt\_swe\_t1\_real\_reuse & 0.000 & patch\_apply\_failed & No & phase107\_gpt\_swe\_t1\_source\_context\_followup & 0.000 & Yes & No & test\_command\_failed & Source context fixed patch application; hidden test still failed \\
\bottomrule
\end{tabular}%
}
\end{table*}

\begin{table*}[t]
\centering
\caption{SWE-T1 issue-aligned follow-up. The first-pass SWE-T1 row remains the main-table evidence; this diagnostic uses the pre-registered issue-aligned scorer and shows that both Summary and PaperToSkill pass under the revised contract, so it is not evidence of PaperToSkill advantage.}
\label{tab:swe-t1-issue-aligned-followup}
\resizebox{\textwidth}{!}{%
\begin{tabular}{llllllllllll}
\toprule
Task & Condition & Main Run & Main Score & Phase107 Run & Phase107 Score & Phase107 Test & Issue Run & Issue Score & Issue Test & Issue Test Patch & Interpretation \\
\midrule
SWE-T1 & Summary & phase97\_gpt\_swe\_t1\_summary\_retry & 0.000 & phase107\_gpt\_swe\_t1\_source\_context\_followup & 0.000 & No & phase110\_gpt\_swe\_t1\_issue\_aligned\_followup & 1.000 & Yes & No & Issue-aligned scorer passes both conditions; no PaperToSkill advantage \\
SWE-T1 & PaperToSkill & phase97\_gpt\_swe\_t1\_real\_reuse & 0.000 & phase107\_gpt\_swe\_t1\_source\_context\_followup & 0.000 & No & phase110\_gpt\_swe\_t1\_issue\_aligned\_followup & 1.000 & Yes & No & Issue-aligned scorer passes both conditions; no PaperToSkill advantage \\
\bottomrule
\end{tabular}%
}
\end{table*}

\begin{table*}[t]
\centering
\caption{SnapATAC2 executable-artifact follow-up. The main SNAP rows in Table~\ref{tab:real-reuse-main} remain unchanged; this paired diagnostic executes a pre-registered controlled scaffold over the same miniature fixtures to test whether concrete artifacts plus runtime/memory records satisfy the existing scorer.}
\label{tab:snapatac2-executable-followup}
\resizebox{\textwidth}{!}{%
\begin{tabular}{lllllll}
\toprule
Task & Condition & Run & Score & Success & Runtime Seconds & Peak Memory MB \\
\midrule
SNAP-T1 & summary & phase108\_snapatac2\_executable\_artifact\_followup & 1.0 & True & 1.296368 & 26.378289 \\
SNAP-T1 & papertoskill & phase108\_snapatac2\_executable\_artifact\_followup & 1.0 & True & 1.209319 & 26.338634 \\
SNAP-T2 & summary & phase108\_snapatac2\_executable\_artifact\_followup & 1.0 & True & 1.75799 & 24.648565 \\
SNAP-T2 & papertoskill & phase108\_snapatac2\_executable\_artifact\_followup & 1.0 & True & 1.756513 & 24.758577 \\
\bottomrule
\end{tabular}%
}
\end{table*}

\begin{table*}[t]
\centering
\caption{SnapATAC2 SNAP-T1 executable-candidate follow-up. The main SNAP rows in Table~\ref{tab:real-reuse-main} remain unchanged; this diagnostic executes model-generated phase112 candidate scripts under the pre-registered executable-candidate contract and shows contract closure for both conditions, not PaperToSkill advantage.}
\label{tab:snapatac2-executable-candidate-followup}
\resizebox{\textwidth}{!}{%
\begin{tabular}{lllllll}
\toprule
Task & Condition & Run & Score & Success & Runtime Seconds & Peak Memory MB \\
\midrule
SNAP-T1 & summary & phase112\_gpt\_snapatac2\_executable\_candidate\_t1 & 1.0 & True & 9.362002 & 0.0 \\
SNAP-T1 & papertoskill & phase112\_gpt\_snapatac2\_executable\_candidate\_t1 & 1.0 & True & 2.010309 & 164.615443 \\
\bottomrule
\end{tabular}%
}
\end{table*}

\begin{table*}[t]
\centering
\caption{Auxiliary Full Excerpt sanity check over the pre-registered subset. Scores come from matched real-reuse rows where available; token columns are local whitespace context proxies, not provider billing or output-token costs.}
\label{tab:full-excerpt-sanity}
\resizebox{\textwidth}{!}{%
\begin{tabular}{llllllllll}
\toprule
Task & Paper & Metric & Summary & PaperToSkill & Full Excerpt & Summary Tokens & PaperToSkill Tokens & Full Excerpt Tokens & Status \\
\midrule
AIDE-T1 & AIDE & validation\_score & 0.816 & 0.817 & 0.000 & 121 & 878 & 7366 & Scored (GPT-family) \\
SWE-T1 & SWE-agent & resolved & 0.000 & 0.000 & 0.000 & 89 & 1173 & 42048 & Scored (GPT-family) \\
SNAP-T1 & SnapATAC2 & runtime\_memory\_quality & 0.000 & 0.500 & 0.250 & 72 & 1069 & 10297 & Scored (GPT-family) \\
\bottomrule
\end{tabular}%
}
\end{table*}

\begin{table*}[t]
\centering
\caption{Main deterministic/offline PaperToSkill results. Coverage scores are task-rubric scores; transfer readiness is an offline artifact-readiness metric.}
\label{tab:main-results}
\begin{tabular}{lrrrrrrrr}
\toprule
Paper & Rubric & Skill & Summary & Abstract & $\Delta$Sum. & Transfer & Support & Words \\
\midrule
AI Scientist-v2 & 20/20 & 7.867/9 & 1.733/9 & 1.2/9 & 6.134 & 10/10 & 0.938 & 782 \\
Reflexion & 20/20 & 8.267/9 & 3.483/9 & 2.533/9 & 4.784 & 10/10 & 1.000 & 479 \\
AIDE & 20/20 & 9.1/10 & 1.916/10 & 1.333/10 & 7.184 & 10/10 & 1.000 & 927 \\
Toolformer & 20/20 & 8.9/10 & 2.5/10 & 1.534/10 & 6.400 & 10/10 & 1.000 & 943 \\
\bottomrule
\end{tabular}
\end{table*}

\begin{table}[t]
\centering
\caption{Transfer-note ablation. Removing only \texttt{Transfer Notes} lowers offline readiness across all curated skills.}
\label{tab:transfer-ablation}
\begin{tabular}{lrr}
\toprule
Paper & Full & No Transfer \\
\midrule
AI Scientist-v2 & 10/10 & 7.6/10 \\
Reflexion & 10/10 & 7.6/10 \\
AIDE & 10/10 & 7.6/10 \\
Toolformer & 10/10 & 7.6/10 \\
\bottomrule
\end{tabular}
\end{table}

\begin{table*}[t]
\centering
\caption{Local tokenizer-aware context-size proxy using the \texttt{o200k\_base} tokenizer. Counts are not provider bills, output-token costs, or success-per-dollar measurements.}
\label{tab:cost-proxy}
\begin{tabular}{lrrrr}
\toprule
Paper & Full Paper Tokens & Skill Tokens & Skill/Full & Reduction \\
\midrule
AI Scientist-v2 & 45212 & 1079 & 0.0239 & 97.61\% \\
Reflexion & 16414 & 703 & 0.0428 & 95.72\% \\
AIDE & 13312 & 1285 & 0.0965 & 90.35\% \\
Toolformer & 20365 & 1255 & 0.0616 & 93.84\% \\
\bottomrule
\end{tabular}
\end{table*}

\begin{table*}[t]
\centering
\caption{Automatic extracted-text note scaffolds are evaluated separately from curated-note cases.}
\label{tab:auto-note}
\begin{tabular}{llrrrrr}
\toprule
Paper & Input & Rubric & Coverage & Transfer & Support & Words \\
\midrule
Toolformer & Curated note & 20/20 & 8.9/10 & 10/10 & 1.000 & 943 \\
Toolformer & Auto note scaffold & 20/20 & 9.3/10 & 10/10 & 1.000 & 1179 \\
AIDE & Curated note & 20/20 & 9.1/10 & 10/10 & 1.000 & 927 \\
AIDE & Auto note scaffold & 20/20 & 8.467/10 & 9.5/10 & 1.000 & 998 \\
\bottomrule
\end{tabular}
\end{table*}
